\section{Related Work}
\label{sec:related}
\paragraph{NBV and active scene acquisition.}
Bircher et al. select exploration motions through a receding-horizon NBV planner
whose branches are valued by newly observable space \cite{bircher2016nbv}.
Surface Edge Explorer instead uses the density and boundaries of existing 3D
observations to direct further surface acquisition \cite{border2018see}.
These approaches illustrate different representations for deciding where
additional measurements are useful. Our work does not argue that generic
exploration objectives are incorrect for mapping. It asks whether their
spatial weighting is appropriate when the endpoint is another robot's
manipulation task. Accordingly, our Generic comparator is an objective-level
control in the same implemented planner, not a claimed reproduction or
head-to-head evaluation of these complete published systems.

\paragraph{Task-oriented perception and active grasping.}
Gualtieri and Platt study how viewpoint selection affects the localization and
confidence of grasp detections \cite{gualtieri2017viewpoint}. Breyer et al.
combine exploration of occluded target regions with continuously updated grasp
predictions and a decision to execute or continue observing
\cite{breyer2022target}. These works establish that the value of a measurement
can depend on a grasping objective rather than only on map completeness.
Our formulation transfers this reasoning across heterogeneous robots: aerial
observations are weighted by the ground mobile manipulator's feasible stance
support, including its extended base footprint. The relevance of a ground cell
is thus derived from future base placements, not solely from whether it belongs
to an incompletely observed target. We do not claim to introduce task-oriented
perception in general.

\paragraph{Aerial-ground cooperation.}
ColAG couples aerial perception and mapping to the navigation needs of
perception-limited ground vehicles, including aerial waypoint scheduling
\cite{li2023colag}. Lenz et al. demonstrate autonomous perception and brick
handling with ground and aerial systems for wall construction
\cite{lenz2020wall}. These studies motivate treating cooperative sensing and
physical action as an integrated robotics problem. Our emphasis is the
intermediate representation linking a ground arm's manipulation support to
aerial viewpoint value, and its evaluation against a matched generic objective
through physical retrieval. We do not claim general multi-robot task allocation
or a replacement for collaborative navigation and localization.

\paragraph{Reachability and mobile manipulation.}
Capability-map work represents the structure of a robot's reachable workspace
\cite{zacharias2007capability}. RM4D provides a compact representation supporting
both reachability and inverse-reachability queries, including mobile-manipulator
base-placement use cases \cite{rudorfer2025rm4d}. We reuse RM4D as the source of
candidate manipulation support rather than presenting it as a contribution of
this paper. The additional step is to project validated base-pose support into
the environment that must be sensed before those poses can be used. Reachability
alone remains insufficient: base occupancy, measured ground support, close-range
visibility and collision-aware execution are evaluated downstream. Our shared
execution interface makes these limitations visible in the experimental outcome
instead of assuming that a valid inverse-kinematics solution completes the task.
